
\documentclass{article}
\usepackage{colortbl}
\usepackage{makecell}
\usepackage{multirow}
\usepackage{supertabular}

\begin{document}

\newcounter{utterance}

\centering \large Interaction Transcript for game `ifeval', experiment `ifeval', episode 158 with reflexion{-}Qwen3.5{-}27B.
\vspace{24pt}

{ \footnotesize  \setcounter{utterance}{1}
\setlength{\tabcolsep}{0pt}
\begin{supertabular}{c@{$\;$}|p{.15\linewidth}@{}p{.15\linewidth}p{.15\linewidth}p{.15\linewidth}p{.15\linewidth}p{.15\linewidth}}
    \# & \multicolumn{2}{c}{Player} && \multicolumn{2}{c}{Game Master} \\
    \hline

    \theutterance \stepcounter{utterance}  
    & & & \multicolumn{4}{p{0.6\linewidth}}{
        \cellcolor[rgb]{0.9,0.9,0.9}{
            \makecell[{{p{\linewidth}}}]{
                \texttt{\tiny{[P1$\langle$GM]}}
                \texttt{You have attempted this same task before and did not succeed. These are the plans you wrote after those attempts.} \\
\texttt{Trial \#1: New plan: I failed because I did not strictly adhere to the formatting constraint of separating the paragraphs with `***`. Although I included the required keywords and wrote three paragraphs, the absence of the specific separator caused the task to fail. Next time, I will ensure that after writing the first paragraph containing "calculate", I insert `***` on its own line, then write the second paragraph containing "file", insert `***` again, and finally write the third paragraph containing "conclusion". I will double{-}check that the separator appears exactly twice between the three distinct blocks of text.} \\
\\ 
\texttt{Trial \#2: New plan: I failed because I did not include the required `***` separator between the paragraphs, even though I used the correct keywords and wrote three distinct blocks of text. The task explicitly required separating the paragraphs with `***`. Next time, I will write the first paragraph containing "calculate", immediately follow it with a line containing only `***`, write the second paragraph containing "file", follow that with another `***`, and then write the final paragraph containing "conclusion". I will verify that the output contains exactly two instances of `***` acting as delimiters between the three required paragraphs.} \\
\\ 
\\ 
\texttt{I was hoping you could help me out with a few things. I wish to learn about how to pay my taxes. Can you summarize the process for me? Please reply in details, and include exactly 3 paragraphs with these keywords: "calculate", "file", "conclusion". Separate the paragraphs with ***.} \\
            }
        }
    }
     \\ \\

    \theutterance \stepcounter{utterance}  
    & \multicolumn{4}{p{0.6\linewidth}}{
        \cellcolor[rgb]{0.561,0.737,0.561}{
            \makecell[{{p{\linewidth}}}]{
                \texttt{\tiny{[P1$\rangle$GM]}}
                \texttt{To begin the process of paying your taxes, you must first gather all necessary financial documents, such as W{-}2s and 1099s, to accurately **calculate** your total taxable income and determine the correct amount of tax you owe or the refund you are due.} \\
\\ 
\texttt{***} \\
\\ 
\texttt{Once you have determined your tax liability, the next step is to **file** your tax return with the appropriate government agency, which can often be done electronically through official tax software or by mailing a completed paper form before the annual deadline.} \\
\\ 
\texttt{***} \\
\\ 
\texttt{In **conclusion**, staying organized with your records and understanding the specific deadlines and requirements for your situation will ensure a smooth tax season and help you avoid potential penalties or interest charges.} \\
            }
        }
    }
    & & \\ \\

    \theutterance \stepcounter{utterance}  
    & & & \multicolumn{2}{p{0.3\linewidth}}{
        \cellcolor[rgb]{0.9,0.9,0.9}{
            \makecell[{{p{\linewidth}}}]{
                \texttt{\tiny{[GM$|$GM]}}
                {'keywords:existence': {'keywords': ['calculate', 'file', 'conclusion']}, 'length_constraints:number_paragraphs': {'num_paragraphs': 3}}
            }
        }
    }
    & & \\ \\

    \theutterance \stepcounter{utterance}  
    & & & \multicolumn{2}{p{0.3\linewidth}}{
        \cellcolor[rgb]{0.9,0.9,0.9}{
            \makecell[{{p{\linewidth}}}]{
                \texttt{\tiny{[GM$|$GM]}}
                \texttt{game\_result = WIN} \\
            }
        }
    }
    & & \\ \\

\end{supertabular}
}

\end{document}
